% =====================================================================
% The 15:50 Candle -- journal version.
% The full working paper in ../spx_candle_1550/ is its Internet Appendix.
% The bibliography is inline (no bibtex pass); Figure 1 reads figures/*.csv.
% Set \blindtrue for a copy without the author's name (double-blind review).
% =====================================================================
\documentclass[12pt,letterpaper]{article}

\newif\ifblind
\blindfalse

\usepackage[utf8]{inputenc}
\usepackage[T1]{fontenc}
\usepackage{amsmath,amsthm}
\usepackage{lmodern}
\usepackage[margin=1in]{geometry}
\usepackage{setspace}
\usepackage{booktabs}
\usepackage{array}
\newcolumntype{L}[1]{>{\raggedright\arraybackslash}p{#1}}
\usepackage{microtype}
\usepackage{natbib}
\usepackage{multirow}
\usepackage{float}
\usepackage{xcolor}
\usepackage{etoolbox}
\usepackage{enumitem}
\usepackage[font=small,labelsep=colon,skip=4pt]{caption}
\captionsetup[table]{position=top}
\captionsetup[figure]{position=top}
\usepackage[compact]{titlesec}
\usepackage{pgfplots}
\pgfplotsset{compat=1.18}
\definecolor{linkblue}{HTML}{1C4E8C}
\usepackage[colorlinks=true,allcolors=linkblue,hyperfootnotes=false]{hyperref}
\ifblind
\hypersetup{pdftitle={The 15:50 Candle: A Pre-Registered Study of an End-of-Day Signal in 0DTE SPX Options}}
\else
\hypersetup{pdftitle={The 15:50 Candle: A Pre-Registered Study of an End-of-Day Signal in 0DTE SPX Options}, pdfauthor={Kausthubh Veldanda}}
\fi

\pagestyle{plain}
\usepackage{titling}

% compact title block: title near the top margin, author, date and abstract close below
\setlength{\droptitle}{-5.5em}
\pretitle{\begin{center}\Large}
\posttitle{\par\end{center}\vspace{-1.3em}}
\preauthor{\begin{center}\normalsize}
\postauthor{\par\end{center}\vspace{-1.5em}}
\predate{\begin{center}\small}
\postdate{\par\end{center}\vspace{-2.4em}}

\setlist{itemsep=1pt,topsep=3pt}

\setcounter{topnumber}{3}
\setcounter{bottomnumber}{2}
\setcounter{totalnumber}{5}
\renewcommand{\topfraction}{0.9}
\renewcommand{\bottomfraction}{0.8}
\renewcommand{\textfraction}{0.07}
\renewcommand{\floatpagefraction}{0.8}
\setlength{\textfloatsep}{10pt plus 4pt minus 3pt}
\setlength{\floatsep}{8pt plus 2pt minus 2pt}
\makeatletter
\setlength{\@fptop}{0pt}
\setlength{\@fpsep}{16pt}
\setlength{\@fpbot}{0pt plus 1fil}
\makeatother
\setlength{\parskip}{0pt plus 1pt}
\flushbottom
\setlength{\intextsep}{8pt plus 4pt minus 2pt}
% table body size; one wide table sets \tabfont to 9pt locally
\newcommand{\tabfont}{\footnotesize}
\AtBeginEnvironment{tabular}{\tabfont}
\AtBeginEnvironment{table}{\setstretch{1}}
\AtBeginEnvironment{figure}{\setstretch{1}}
\setlength{\emergencystretch}{2em}

% the explanatory note of a table or figure, set below it in small type
\newcommand{\tabnotes}[1]{\par\vspace{5pt}\parbox{\linewidth}{\footnotesize\setlength{\parindent}{0pt}\textit{Note:} #1}\par}

\definecolor{series1}{HTML}{2A78D6}
\definecolor{series2}{HTML}{EB6834}
\pgfplotsset{
  table/col sep=comma,
  every axis/.append style={font=\fontsize{9}{11}\selectfont},
  paper/.style={
    axis lines*=left, axis line style={black!60}, tick style={black!60},
    tick align=outside, major tick length=2pt,
    ymajorgrids, grid style={black!10},
    label style={font=\fontsize{10}{12}\selectfont},
    legend style={draw=none, fill=none, font=\fontsize{9}{11}\selectfont, legend cell align=left},
    /pgf/number format/1000 sep={,}, scaled y ticks=false,
    y tick label style={/pgf/number format/fixed},
  },
  zeroline/.style={extra y ticks={0}, extra y tick labels={},
    extra y tick style={grid=major, major grid style={black!55, thin}, tick style={draw=none}}},
}

\newcommand{\ITM}[1]{\textsc{itm}#1}
\newcommand{\OTM}[1]{\textsc{otm}#1}
\newcommand{\ATM}{\textsc{atm}}
\newcommand{\dollar}{\$}

\begin{document}
\ifblind
\title{\textbf{The 15:50 Candle: A Pre-Registered Study\\of an End-of-Day Signal in 0DTE SPX Options}\thanks{This version: 12 November 2026.
Every confirmatory test was pre-registered: each plan was hashed and time-stamped
before the data it governs was read, and the tests that failed are reported with
those that passed. An Internet Appendix reports every study in full. A statement
on data availability is at the end of the paper. This paper is research, not
investment advice.}}
\author{Anonymous for review}
\else
\title{\textbf{The 15:50 Candle: A Pre-Registered Study\\of an End-of-Day Signal in 0DTE SPX Options}\thanks{This version: 12 November 2026.
Every confirmatory test was pre-registered: each plan was hashed and time-stamped
before the data it governs was read, and the tests that failed are reported with
those that passed. An Internet Appendix, the long working-paper version of this
study, reports every study in full. A statement on data availability is at the
end of the paper. I thank Tiger Chen, whose observation started this study.
This paper is research, not investment advice.}}
\author{Kausthubh Veldanda\thanks{University of California, Berkeley. Email: \href{mailto:kveldanda@berkeley.edu}{kveldanda@berkeley.edu}}}
\fi
\date{12 November 2026}
\maketitle

\begin{abstract}
\noindent I test whether a one-minute candle near the close predicts the move of the
S\&P~500 index (SPX) into the close, and whether the prediction is worth
anything after the cost of trading same-day-expiry (0DTE) index options. Every confirmatory test was
pre-registered. A scan of 41 candle features on 2013--2020 selected the 15:50 candle, read from a CFD that tracks the index, when its body
spans at least 60\% of its range. On 2021--2026 this signal was right on 57.5\%
of 644 trades, and an SPXW option 10 points in the money, bought at the quoted
ask and held to settlement, earned \$67 a trade ($t = 2.40$). The profit
depends on the bid-ask spread. The same trades lost money before 2021, when spreads
were about five times wider, and limit orders at the mid did not help because they
missed the best days. Twenty-six pre-registered stop and take-profit rules, tested
at real bids, all earned less than holding on 2021--2023. The result also depends on where the candle is read. Read from the index
itself, the
rule missed its criterion ($t = 1.53$). In a test using the exchanges'
imbalance messages, the candle
pointed the same way as the closing-auction imbalance published at 15:50 on 73\%
of signal days but was not more accurate on those days (56\% against 62\%).

\medskip
\noindent \textit{JEL:} G13, G14\\
\noindent \textit{Keywords:} 0DTE options; S\&P~500 index options; closing
auction; intraday predictability; bid-ask spread; pre-registration.
\end{abstract}



% =====================================================================
\section{Introduction}
\label{sec:intro}

At 15:50~ET, ten minutes before the U.S. equity market closes, the New York
Stock Exchange publishes its regulatory closing imbalance and Nasdaq begins to
disseminate an early imbalance indicator for the close. From then on the NYSE
accepts new market-on-close and limit-on-close orders only on the side that
offsets the published imbalance.\footnote{See the NYSE fact sheet \emph{NYSE Opening and
Closing Auctions} and Nasdaq's \emph{Closing Cross} FAQ. The NYSE moved its
cut-off and imbalance publication from 15:45 to 15:50 on 1 April 2019, under a rule
change the SEC approved on 31 January 2019 (Release No.~34-85021). Nasdaq moved
the other way: in the fourth quarter of 2018 its cut-off and the first
dissemination of its Net Order Imbalance Indicator moved from 15:50 to 15:55
(Release No.~34-84454), and it later added an abbreviated Early Order Imbalance
Indicator from 15:50, scheduled for the fourth quarter of 2019 (Release
No.~34-87179).} An option on the
S\&P~500 index that expires that day and is bought at 15:51 settles in cash on
the official close nine minutes later, so its payoff depends on little more than
the direction and size of the last move of the session. Before 2016 such
same-day (0DTE) SPXW expiries were listed only on some days; since 2022 they have
been listed daily. If the first minute after the imbalance is published says anything about where
the index will close, these options are the simplest way to trade on it.

I came to this question through a friend\ifblind\else, Tiger Chen\fi, who had
noticed by chance, on one-minute charts of the index, that the S\&P~500 often
makes a sharp move in one direction about ten minutes before the close, and who
showed me 27 screenshots of it. This paper began as an attempt to check whether
that observation is real and whether trading it with same-day options makes
money.

A backtest of such a trade can mislead in several ways. If enough candle rules
are scored on enough history, one of them will look good by chance, and its
edge will shrink out of sample \citep{harvey2016}. A backtest that enters at
the close of the signal candle uses information that was not yet available. And
a rule that is right 54\% of the time can still lose money once the bid-ask
spread is paid. I therefore fixed what would be tested before reading the data,
and priced the trade at quotes that were available at the time.

Every confirmatory test in the paper was pre-registered \citep{nosek2018}.
Before reading the data for a study I wrote down what would be computed, on
which data, and what would count as success, and the plan was hashed and
time-stamped. Every study is reported, including the ones that failed. The work
went in four steps. I first scanned 41 candle features on 2013--2020 and tested
the ones that qualified on 2021--2026 index moves. I then priced the trade with
real OPRA quotes, tuning on 2013--2020 and confirming once on 2021--2026.
Third, I used trade bars and quote snapshots to estimate what execution costs.
Fourth, because the backtests read the signal from a contract for difference
(CFD) that tracks the index, I read the same candle from the SPX index and from
SPY, and searched the other closing minutes for a better candle. A last
pre-registered test uses the exchanges' imbalance messages to ask whether the
candle is more accurate when it agrees with the imbalance published at 15:50.

The scan produced one usable signal, the 15:50 one-minute candle with a
\emph{strong body}, traded in its own direction. On 2021--2026 it was right on
57.5\% of 644 trades. It is a stricter version of an effect I had already seen
in earlier exploratory work, and it has weakened since 2023.

At real option prices the edge is small. A deep
in-the-money option bought at the quoted ask and held to settlement earned
$+\$67.46$ per contract on 2021--2026 ($t = 2.40$), about two-thirds of what a
calibrated pricing model implied. Whether the trade makes money depends mostly on the bid-ask spread, and the
same trades lost money at the ask before 2021, when spreads were wider.
A limit order at the mid does not recover the spread, because it fills mainly
when the trade is going wrong, and a one-minute delay costs a sixth of the edge. Selling before the close did not help either. All 26 pre-registered exit rules
that watch the option's bid earned less than holding on 2021--2023.

The result also depends on where the candle is read. Under a plan locked before
any index candle was joined with a price, the rule read from the SPX index
candle made $+\$39.44$ a trade ($t = 1.53$) and missed its criterion. Read from
the SPY candle it made $+\$71.74$ ($t = 2.35$) and met a retrospective
criterion. On all three sources the gain came from sessions on which every
candle showed a strong body.

I did not find a better rule. A search over 7{,}128 settings of
candle source, minute, cut, direction and strike chose the 15:50 candle again,
although its choice was not clearly beyond luck ($p = 0.051$). The timing of the closing auction is the obvious reason this minute might
matter, and the exchanges' imbalance data allow one direct check. The candle
pointed the same way as the closing-auction imbalance
published at 15:50 on 73\% of signal days, but in a pre-registered test it was not
more accurate on those days (56\% against 62\%). The auction explanation was therefore not confirmed, although the test could detect only a large
difference.

Several limits apply throughout, and none of this is a recommendation to trade.
The holdout was not blind, because a related effect was already known. The
real-price confirmation reuses index outcomes that I had already seen, and the
execution results are descriptive proxies. The confirmed rule reads its candle
from a CFD and not from the index, and the SPY candle's pass is retrospective.

\paragraph{Related literature.} The closest work is on market intraday momentum. \citet{ghlz2018} find
that the first half-hour return predicts the last half hour of the S\&P~500, and
\citet{bdlm2021} find the same pattern across futures markets and link it to the
hedging of option dealers who are short gamma. The 15:50 candle lies inside that
last half hour, but in my data neither momentum nor a dealer-gamma proxy explains it
or improves it. What I add is a single minute within that window, tested out of
sample under a pre-registered plan. On the closing auction, \citet{bm2023} show that closing-auction volume is tied to
passive and ETF ownership and that auction imbalance information contributes to
price discovery at the close. My signal sits in the first minute after that information is published, and I
use the exchanges' imbalance messages to test whether its accuracy depends on
agreeing with it. On 0DTE options, \citet{bbg2023} document that retail
traders lose money in them, and \citet{dev2024} find little evidence that 0DTE
gamma or trading volume amplifies index moves. I look at a narrower question, which is how much of a short-horizon prediction
is left once the spread on these contracts is paid. Finally, the paper follows
the work on data mining and pre-registration in return predictability
\citep{harvey2016,nosek2018}. It applies a plan, hash and ledger protocol to a
sequence of studies and reports the tests that failed.

\paragraph{Outline.} Section~\ref{sec:data} describes the data and the design.
Section~\ref{sec:scan} tests the signal, and Sections~\ref{sec:realprice}
and~\ref{sec:costs} price the trade at real quotes and study its costs.
Section~\ref{sec:source} examines the source of the candle and
Section~\ref{sec:search} the search for a better rule. Section~\ref{sec:why}
discusses why 15:50 might matter and tests one explanation, and
Section~\ref{sec:conclusion} concludes.

% =====================================================================
\section{Data and research design}
\label{sec:data}



\subsection{Data}

Table~\ref{tab:data} lists the data sources.

\begin{table}[htbp]
\centering
\caption{Data sources.}
\label{tab:data}
\begin{tabular}{@{}L{0.17\linewidth}L{0.41\linewidth}L{0.13\linewidth}L{0.21\linewidth}@{}}
\toprule
Source & Content & Period & Used for \\
\midrule
Dukascopy & US500 CFD one-minute bars, regular trading hours & 2013--2026 & Signals; index outcomes \\
Databento OPRA & SPXW best bid and offer at one-minute marks, 15:00--16:01~ET, the 25 strikes nearest the money, calls and puts & 2013--2026 & Option prices \\
Databento OPRA & SPXW one-minute trade bars for the same contracts & 2013--2026 & Execution \\
Databento OPRA & Daily open interest, all SPX and SPXW options & 2013--2026 & Gamma proxy \\
Databento OPRA & XSP and SPY quotes & 2021--2026 & Smaller contracts \\
Cboe & Official SPX daily close & 2013--2026 & Option settlement \\
CNBC & SPX index one-minute bars (archive) & 2021--2026 & Index candle (Section~\ref{sec:source}) \\
CNBC & SPY one-minute bars (archive; traded prices, not dividend-adjusted) & 2021--2026 & SPY candle (Section~\ref{sec:source}) \\
Databento & NYSE and Nasdaq closing-auction imbalance messages, 15:50--16:00~ET on signal days, for SPY's 50 largest holdings (SEC N-PORT reports) & 2021--2026 & Imbalance test (Section~\ref{sec:imbtest}) \\
\bottomrule
\end{tabular}
\tabnotes{The index sample runs from 2013-02-28 to 2026-10-20 and the OPRA SPXW
sample from 2013-04-01 to 2026-10-20. Robinhood one-minute SPX and SPXW bars,
used to calibrate the pricing model and for the forward record, begin on
2026-09-21.}
\end{table}

Signals and index outcomes come from Dukascopy's US500 CFD one-minute bars during
regular trading hours, restricted to a validated set of full sessions. Discovery
uses 2013--2020 (1{,}861 sessions) and the holdout 2021--2026 (1{,}446 sessions).
The CFD stands in for SPX, and index outcomes are measured to the close of the
15:59 bar rather than to the official settlement. Section~\ref{sec:source} reads
the signal candle instead from archives of SPX and SPY one-minute bars.

Real option prices come from Databento's OPRA feed: the consolidated best bid and
offer at one-minute marks for SPXW contracts expiring that day, from 15:00 to
16:01~ET, for the 25 strikes within 60 points of the at-the-money strike at
15:51. A quote stamped HH:MM:00 is the quote at the open of that minute.
One-minute trade bars for the same contracts support the execution analysis, and
daily open interest on every SPX and SPXW option supports a dealer-gamma proxy.
Of the 3{,}287 index days, 2{,}149 had an SPXW expiry that day; the rest are
excluded from every option study. Options settle at intrinsic value on Cboe's
official SPX close.

\subsection{The signal and the trade}
\label{sec:signal}

Let $o, h, l, c$ be the open, high, low and close of a one-minute bar. The bar
has a \emph{strong body} if $h > l$ and
\begin{equation}
|c - o| \ge 0.6\,(h - l),
\end{equation}
and its direction is the sign of $c - o$. The signal is a strong body on the
15:50 bar. Because that bar is complete only at 15:50:59, the trade enters at the
open of the 15:51 bar and, in the index tests, is scored on the move to the close
of the 15:59 bar.

In the option studies, let $S$ be the index level at entry and
$K_{\textsc{atm}} = 5\lfloor S/5 + 1/2\rfloor$ the nearest listed strike. An
offset of $j$ points, positive meaning in the money in the trade's direction,
buys the call struck at $K_{\textsc{atm}} - j$ after an up signal and the put
struck at $K_{\textsc{atm}} + j$ after a down signal. I consider nine offsets,
from 20 points in the money (\ITM{20}) to 20 points out of the money (\OTM{20})
in steps of five. One contract (multiplier 100) is bought per signal and held to
cash settlement, and each side pays a fee of \$0.09 per contract, with \$0.74 as
a sensitivity. I refer to the 15:50 strong-body signal traded at \ITM{10} as the \emph{rule}.

\subsection{Pre-registration}
\label{sec:prereg}

Each confirmatory study was specified in a written plan before the data it
governs was read. The plan's SHA-256 hash and a UTC timestamp were appended to a
ledger, and each study's code recomputes the hash at start-up and refuses to read
real data unless it matches a ledger entry. Each plan lists what was already
known when it was written, and any departure from a plan is reported as a
labelled deviation; no result was re-run under a changed rule.
Table~\ref{tab:studies} lists every study, its primary criterion and its outcome.
Analyses that were not pre-registered are labelled \emph{descriptive} and were
not used to decide any test.

\begin{table}[htbp]
\centering
\caption{The pre-registered studies.}
\label{tab:studies}
\begin{tabular}{@{}L{0.16\linewidth}L{0.07\linewidth}L{0.11\linewidth}L{0.33\linewidth}L{0.22\linewidth}@{}}
\toprule
Study & Section & Registered & Primary criterion & Outcome \\
\midrule
Candle scan & \ref{sec:scan} & 2026-10-21 & Holm-corrected test of each frozen candidate on 2021--2026, and an accuracy bar & One of three candidates passes \\
Strike ladder & \ref{sec:strikes} & 2026-11-08 & Nine strikes ranked under a calibrated model & Deep in the money most robust \\
Early exits & \ref{sec:other} & 2026-11-09 & A rule beats holding with $t > 2.64$, robust across model variants & No rule passes \\
Real option prices & \ref{sec:realprice} & 2026-11-09 & Switch only if a new setting leads by $t \ge 2.0$; confirm with $t \ge 2.0$ & Rule kept and confirmed \\
Re-tuning & \ref{sec:retune} & 2026-11-09 & Luck diagnostic $p \le 0.05$, then $t \ge 2.0$ on 2024--2026 & Pick fails the luck diagnostic \\
Intraday momentum & \ref{sec:other} & 2026-11-09 & Slope $t \ge 2.0$; filters $p \le 0.0125$; gamma $p \le 0.025$ & Nothing supported \\
XSP and SPY & \ref{sec:other} & 2026-11-09 & Positive mean with $t \ge 2.0$ on 2021--2026 & Neither confirmed \\
Index candle & \ref{sec:source} & 2026-11-10 & Rule read from the index candle: positive mean per day with $t \ge 2.0$ & Criterion not met \\
SPY candle & \ref{sec:source} & 2026-11-10 & Rule read from the SPY candle: positive mean per day with $t \ge 2.0$ & Retrospective criterion met \\
Candle search & \ref{sec:minutes} & 2026-11-10 & Best of 7{,}128 settings on 2021--2023, tested on 2024--2026 only if not at 15:50 and luck diagnostic $p \le 0.05$ & Pick at 15:50; no test \\
Closing imbalance & \ref{sec:imbtest} & 2026-11-10 & Hit rate higher on days the candle agrees with the 15:50 imbalance: $d > 0$ with $t \ge 2.0$ & Not confirmed \\
Real-time exits & \ref{sec:exits} & 2026-11-10 & An exit rule beats holding in both halves of 2021--2023 and has luck diagnostic $p \le 0.05$; then $t \ge 2.0$ on 2024--2026 & Stopped at the coverage gate \\
Real-time exits, revised & \ref{sec:exits} & 2026-11-11 & The same, for the 26 rules that read only the option's bid & No eligible rule; no test \\
Forward options log & \ref{sec:other} & 2026-10-30 & Running record of real option outcomes & Four priced trades so far \\
\bottomrule
\end{tabular}
\tabnotes{Registration dates are UTC. The strike-ladder and early-exit studies
priced trades with a calibrated model, before real quotes were available; they
and the forward log are reported in full in the Internet Appendix. The
limit-order and fixed-minute sale analyses of Section~\ref{sec:costs} were not
pre-registered.}
\end{table}

\subsection{Statistical inference}
\label{sec:inference}

I report mean P\&L per trade and per day, where the daily series takes the
trade's P\&L on a signal day and zero otherwise. Unless stated otherwise,
$t$-statistics use Newey--West standard errors with five lags on the daily series
\citep{neweywest1987}. The scan tests its candidates one-sided with a Holm
correction \citep{holm1979}, and hit-rate intervals are Wilson intervals
\citep{wilson1927}.

Studies that choose one setting from a large grid also report a \emph{luck
diagnostic}. The signal series is shifted circularly against the outcome days,
by every shift from 21 days to $n - 21$ in steps of three. A shift keeps the
signal's own structure but breaks its link to the outcomes. The whole search is
repeated for each shift and its best score kept. With $S$ shifts, of which $k$
produce a best score at least as good as the real one, $p = (1 + k)/(1 + S)$.
The diagnostic asks whether the real winner beats the winner that unrelated
signals would find. It is approximate, because expiry frequency, volatility and
pricing change over the sample.

% =====================================================================
\section{Does the 15:50 candle predict the close?}
\label{sec:scan}



The scan defined 41 features, each of which outputs up, down or no trade using
only bars completed before the entry: the direction of every one-minute candle
from 15:30 to 15:50; five candle shapes (strong body, hammer, shooting star,
engulfing, big candle) on the 15:50 one-minute candle and on the 15:45--15:49
five-minute candle; five-minute directions; and four whole-day features. Features
with $|t| \ge 2.5$ on 2013--2020 qualified, up to eight. The list of qualifying
features was frozen and hashed before any 2021 bar was read. On 2021--2026 each
candidate faced a one-sided test with a Holm correction at 5\%. A candidate that
passed was called \emph{accurate} only if the lower bound of its 95\% Wilson
interval exceeded the hit rate at which buying the at-the-money option breaks
even.

\begin{table}[htbp]
\centering
\caption{The candle scan: frozen candidates on the 2021--2026 holdout.}
\label{tab:scan}
\begin{minipage}{\linewidth}\centering
\textit{Panel A. The three candidates}\par\smallskip
\setlength{\tabcolsep}{4pt}
\begin{tabular}{lcccccc}
\toprule
Candidate & 2013--20 & 2021--26 (95\% CI) & trades & Holm $p$ & predicts & accurate \\
\midrule
\textbf{15:50 strong body} & 54.1\% & \textbf{57.5\%} (53.6--61.2) & 644 & $6.3\times10^{-5}$ & yes & \textbf{yes, just} \\
15:50 direction, any body & 52.2\% & 54.0\% (51.4--56.6) & 1{,}445 & 0.011 & yes & no \\
Day so far (09:30--15:49) & 52.0\% & 48.4\% (45.8--51.0) & 1{,}446 & 0.53 & no & no \\
\midrule
\emph{P3 reference (not tested)} & \emph{51.8\%} & \emph{53.6\%} (51.1--56.2) & \emph{1{,}442} & --- & --- & --- \\
\bottomrule
\end{tabular}
\end{minipage}

\bigskip
\begin{minipage}{\linewidth}\centering
\textit{Panel B. The strong-body candle by subperiod (descriptive)}\par\smallskip
\begin{tabular}{lcccc}
\toprule
Strong body & trades & hit rate (95\% CI) & $t$ & option P\&L / trade ($t$) \\
\midrule
2021--2022 & 216 & 63.0\% (56.3--69.1) & 3.56 & $+\dollar80.7$ (2.33) \\
2023--2026 & 428 & 54.7\% (49.9--59.3) & 2.27 & $+\dollar28.6$ (1.11) \\
\bottomrule
\end{tabular}
\end{minipage}
\tabnotes{Panel A gives hit rates in discovery (2013--2020) and on the holdout,
with 95\% Wilson intervals. The strong-body break-even is 49.5\% at a \$0.09 fee
(51.8\% size-neutral); with 644 trades, ``accurate'' requires an observed hit
rate of 53.4\% (55.7\% size-neutral). The P3 row is a reference and was not
tested. Panel B prices the at-the-money option with the scan's cost model. By
year from 2021 to 2026 the hit rate was 64.9, 61.0, 54.0, 53.8, 57.3 and
53.3\%.}
\end{table}

Panel A of Table~\ref{tab:scan} shows that three features qualified. Two of them,
both on the 15:50 candle, predicted the 2021--2026 move, but only the strong-body
candle also cleared the accuracy bar. It was right on 57.5\% of 644 trades, and
its trades moved an average of 1.82 index points in the signal's direction over
nine minutes, at today's index level. The third candidate, a whole-day feature,
failed out of sample.

I ran two checks, one against chance and one against a generic intraday
pattern. When the
41 features were shifted out of line with the 2013--2020 outcomes and the scan
was re-run 304 times, about 1\% of runs found as many qualifying features as the
real scan, and none found one as strong as the best real feature. The same rules
moved four hours earlier, entering at 11:51, showed nothing in either period.

Both passing
features sit on the 15:50 candle and agree with the direction of the move from
15:50 to 15:51 (an effect I call P3, found in earlier exploratory work) on 98\% and 93\% of
days. I knew how P3 had behaved over 2013--2026 before running the scan, so the
holdout was not blind for these features. I therefore read the result as a
known effect passing a stricter test, not as a new discovery. The strong body
does matter, because on 15:50 candles without one, trading the direction earned
nothing.

The effect has also weakened since 2023. Panel B of Table~\ref{tab:scan} shows
that most of the edge comes from 2021--2022, and in 2023--2026 the hit-rate
interval includes 50\%. At that effect size, a test with 80\% power would need
roughly 19 years of signals, so no single year can settle the question.

% =====================================================================
\section{The trade at real option prices}
\label{sec:realprice}



\subsection{Design}

At the entry minute, the spot $S$ is implied by put--call parity: it is the
median of $K + C_{\text{mid}} - P_{\text{mid}}$ over the three strikes where the
call and put mids are closest. The trade buys at the ask at the open of the entry
minute and settles on the official close, so its P\&L per contract is
\begin{equation}
\text{P\&L} = 100\,\bigl(\text{intrinsic value at the official close} - \text{ask}\bigr) - \text{fee}.
\end{equation}
A trade with no implied spot or no valid quote is dropped and counted.

The study first tuned on 2013--2020 over a grid of 2{,}178 settings: the candle
minute (any one-minute candle from 15:40 to 15:50, entering at the next minute's
open), the body threshold (0.40 to 0.90 in steps of 0.05), the direction (with or
against the candle) and the strike (\ITM{20} to \OTM{20}). A setting's score is
its mean P\&L per day averaged with its neighbours at adjacent thresholds, which
favours a plateau over a sharp peak, and a setting is eligible with at least 150
trades, at least 50 in each half of the period, and a positive mean in both
halves. The best eligible setting would replace the rule only if its luck
diagnostic gave $p \le 0.05$ and its per-day lead over the rule had $t \ge 2.0$.
The chosen rule is confirmed on 2021--2026 if its mean P\&L per day is positive
with $t \ge 2.0$. The plan states that the 2021--2026 index outcomes had already
been used to test the signal, so the confirmation is fresh for option prices but
not for the index moves.

On 2013--2020 spreads were wide and many signal days could not be priced, so
only two of the 2{,}178 settings were eligible. The better one, the 15:42 candle
with a 60\% body at the money, passed the luck diagnostic ($p = 0.034$), but its
lead over the rule was not significant ($t = 1.56$), and the rule went forward
unchanged. The rule itself lost money at the ask in this period.

\subsection{Confirmation on 2021--2026}
\label{sec:confirm}

\begin{table}[htbp]
\centering
\caption{Confirmation: the rule at real asks, 2021--2026.}
\label{tab:rpconf}
\setlength{\tabcolsep}{4pt}
\begin{tabular}{lrrrrrrrr}
\toprule
Period & days & \$/day & $t$ & trades & \$/trade & won & median & worst \\
\midrule
\textbf{2021--2026} & \textbf{1{,}321} & $\mathbf{+29.87}$ & \textbf{2.40} & \textbf{585} & $\mathbf{+67.46}$ & \textbf{54\%} & $\mathbf{+44.91}$ & $\mathbf{-2{,}210}$ \\
2021--2022 & 376 & $+48.21$ & 2.01 & 157 & $+115.46$ & 57\% & $+65.91$ & $-1{,}361$ \\
2023--2026 & 945 & $+22.58$ & 1.56 & 428 & $+49.85$ & 53\% & $+35.91$ & $-2{,}210$ \\
\bottomrule
\end{tabular}
\tabnotes{The rule (15:50 candle, body at least 60\%, with the candle, \ITM{10},
held to the official close) at real asks, \$0.09 fee, dollars per contract. $t$
is Newey--West on the per-day series. Subperiods are descriptive. Mean per trade
by year: 2021 $+\$175$ (67 trades), 2022 $+\$71$ (90), 2023 $+\$4$ (113), 2024
$+\$20$ (106), 2025 $+\$65$ (117), 2026 $+\$121$ (92).}
\end{table}

Table~\ref{tab:rpconf} shows that the rule was confirmed. It earned $+\$29.87$ a
day with $t = 2.40$, against a bar of 2.0. The mean trade made $+\$67.46$,
slightly more than half of all trades were profitable, and the worst lost
\$2{,}210 on one contract. A five-day block bootstrap puts the mean between
$+\$21$ and $+\$115$ a trade (90\% interval), so the mean is positive but
imprecisely measured. Figure~\ref{fig:rpyears} shows the full record.

This confirmation has three weaknesses. It confirms the rule at real prices but
does not show that it is the best setting, since no other setting was evaluated
on this period (Section~\ref{sec:search} returns to that question). I had
already seen the index outcomes. And the edge is weaker since
2023, when it is not significant on its own.

\begin{figure}[htbp]
\centering
\caption{The rule at real option prices.}
\label{fig:rpyears}
\begin{tikzpicture}
\begin{axis}[paper, width=0.96\linewidth, height=5.6cm, xmin=2014, xmax=2027,
  xtick={2014,...,2026}, xticklabel style={/pgf/number format/1000 sep=},
  ymin=-20000, ymax=100000, ytick={-20000,0,...,100000},
  ylabel={Cumulative P\&L (\$)}, zeroline,
  title={(a) Running total, one contract per signal day}, title style={font=\fontsize{10}{12}\selectfont},
  legend style={at={(0.02,0.80)}, anchor=north west}]
\addplot[series1, line width=1.2pt] table[x=t, y=ask_cum]{figures/cumulative.csv};
\addlegendentry{Bought at the ask (as tested)}
\addplot[series2, line width=1pt, dash pattern=on 4pt off 2pt] table[x=t, y=mid_cum]{figures/cumulative.csv};
\addlegendentry{Bought at the mid (an upper bound; see Section~\ref{sec:fills})}
\draw[black!50, densely dotted, thick] (axis cs:2021,-20000) -- (axis cs:2021,100000);
\node[font=\scriptsize, black!70, anchor=north east] at (axis cs:2020.9,98000) {2013--2020: tuning};
\node[font=\scriptsize, black!70, anchor=north west] at (axis cs:2021.1,98000) {2021--2026: one confirmation test};
\end{axis}
\end{tikzpicture}

\vspace{2pt}
\begin{minipage}[t]{0.49\linewidth}\centering
\begin{tikzpicture}
\begin{axis}[paper, ybar, width=0.9\linewidth, height=5.0cm, bar width=6.5pt,
  xmin=2013.4, xmax=2026.6, xtick={2014,2016,...,2026}, minor xtick={2015,2017,...,2025},
  xticklabel style={/pgf/number format/1000 sep=},
  ymin=-8000, ymax=14000, ylabel={P\&L per year (\$)}, zeroline,
  title={(b) Total per calendar year, at the ask}, title style={font=\fontsize{10}{12}\selectfont}]
\addplot[fill=series1, draw=none] table[x=year, y=ask_total]{figures/by_year.csv};
\end{axis}
\end{tikzpicture}
\end{minipage}\hfill
\begin{minipage}[t]{0.49\linewidth}\centering
\begin{tikzpicture}
\begin{axis}[paper, width=0.93\linewidth, height=5.0cm,
  xmin=2013.4, xmax=2026.6, xtick={2014,2016,...,2026}, minor xtick={2015,2017,...,2025},
  xticklabel style={/pgf/number format/1000 sep=},
  ymin=0, ymax=4.5, ytick={0,1,...,4}, ylabel={Median full spread (\$ a share)},
  title={(c) Bid-ask spread of the contract bought}, title style={font=\fontsize{10}{12}\selectfont}]
\addplot[series1, line width=1.2pt, mark=*, mark size=1.6pt] table[x=year, y=median_spread]{figures/by_year.csv};
\end{axis}
\end{tikzpicture}
\end{minipage}
\tabnotes{The rule (15:50 candle, body at least 60\%, with the candle, \ITM{10},
held to the official close) at real Databento OPRA prices, \dollar0.09 fee,
dollars per contract at each year's prices. (a) The 801 priced signal days,
2014--2026. At the ask the rule lost \dollar4{,}997 over the 216 trades of
2013--2020 and made \dollar39{,}463 over the 585 of 2021--2026. The dashed line
buys at the mid and is only an upper bound: a limit order at the mid made about
what paying the ask made, because it missed the best days
(Section~\ref{sec:fills}). (b) 2023 and 2024 were close to zero. (c) Median ask
minus bid at entry. The narrowing spread is why the same signal lost money at the
ask before 2021 and made money after it. Before 2016 SPXW had a same-day expiry
only on some days, so the early years have few trades. Sources: Databento OPRA;
Cboe; Dukascopy.}
\end{figure}

\subsection{Strikes}
\label{sec:strikes}

\begin{table}[htbp]
\centering
\caption{The signal at every strike, real asks, 2021--2026.}
\label{tab:rpstrikes}
\begin{tabular}{lrrrrrrr}
\toprule
Strike & trades & \$/trade & won & median & worst & \$/day & $t$ \\
\midrule
\ITM{20} & 586 & $+77.15$ & 54\% & $+51.91$ & $-2{,}670$ & $+34.20$ & 2.61 \\
\ITM{15} & 584 & $+71.00$ & 55\% & $+52.41$ & $-2{,}330$ & $+31.41$ & 2.43 \\
\textbf{\ITM{10}} & \textbf{585} & $\mathbf{+67.46}$ & \textbf{54\%} & $\mathbf{+44.91}$ & $\mathbf{-2{,}210}$ & $\mathbf{+29.87}$ & \textbf{2.40} \\
\ITM{5} & 585 & $+53.90$ & 50\% & $+2.91$ & $-1{,}900$ & $+23.87$ & 2.07 \\
\ATM & 583 & $+39.50$ & 39\% & $-110.09$ & $-1{,}640$ & $+17.46$ & 1.87 \\
\OTM{5} & 585 & $+12.18$ & 19\% & $-65.09$ & $-1{,}390$ & $+5.40$ & 0.78 \\
\OTM{10} & 581 & $+1.04$ & 11\% & $-25.09$ & $-1{,}180$ & $+0.46$ & 0.10 \\
\OTM{15} & 566 & $-6.96$ & 5\% & $-15.09$ & $-1{,}000$ & $-3.02$ & $-1.03$ \\
\OTM{20} & 529 & $-10.02$ & 3\% & $-10.09$ & $-840$ & $-4.19$ & $-2.34$ \\
\bottomrule
\end{tabular}
\tabnotes{\$0.09 fee, dollars per contract. Descriptive, not a second test:
choosing a strike from this table would be choosing on the confirmation period.}
\end{table}

Before real quotes were available, a pre-registered study priced the holdout
trades with a Bachelier model \citep{bachelier1900}, calibrated on seven days of
real broker prices, and found strikes 10 to 20 points in the money the most
robust. Table~\ref{tab:rpstrikes} shows that at real prices the strikes rank the
same way. Deep in-the-money options carry little time value, so they keep nearly all of a
small move, and \ITM{10} needs about half the premium of \ITM{20}. At the money the mean was positive but most trades lost, and ten or more
points out of the money the option rarely paid.

The model had been too optimistic about the size of the profit. On the same
days and contracts it implied $+\$100.70$ a trade at \ITM{10}, and the real
trade made $+\$67.77$. In the money the model
priced the option well, with the real mid within about 2\% of the model price,
but the real ask sat further above the mid than the model's flat half-spread
allowed.

% =====================================================================
\section{Costs and execution}
\label{sec:costs}



\subsection{The bid-ask spread}
\label{sec:spread}

\begin{table}[htbp]
\centering
\caption{The rule bought at the ask and at the mid.}
\label{tab:spread}
\begin{tabular}{llrrrrrc}
\toprule
Period & bought at & trades & \$/day & $t$ & \$/trade & won & median full spread \\
\midrule
\multirow{2}{*}{2013--2020} & ask & 216 & $-7.70$ & $-0.47$ & $-23.14$ & 42\% & \multirow{2}{*}{\dollar2.90} \\
 & mid & 216 & $+48.19$ & 2.87 & $+144.78$ & 56\% & \\
\midrule
\multirow{2}{*}{2021--2026} & ask & 585 & $+29.87$ & 2.40 & $+67.46$ & 54\% & \multirow{2}{*}{\dollar0.60} \\
 & mid & 585 & $+46.90$ & 3.71 & $+105.90$ & 56\% & \\
\bottomrule
\end{tabular}
\tabnotes{\ITM{10}, \$0.09 fee, dollars per contract. Rows at the mid are
descriptive and decide nothing. The spread is ask minus bid per share at entry,
the median over the priced trades. With a \$0.74 fee the ask rows become
$-\$23.79$ and $+\$66.81$ a trade ($t = 2.37$ in 2021--2026).}
\end{table}

Table~\ref{tab:spread} and Figure~\ref{fig:rpyears}(c) show how much the result depends on the spread. At the mid, the same trades made money in both periods. At the ask they
lost money in 2013--2020, when the median quoted spread was about five times its
2021--2026 level, and made money afterwards. Even since 2021, paying the ask instead of the mid costs about \$39 a trade,
more than a third of the edge at the mid, so how the order is filled matters a
great deal.

The ask is only a quoted price, and I do not observe what an order would
actually have paid. OPRA trade bars give a partial check. Since 2021 the
contract traded in the entry minute on 99\% of
signal days, against just over half the days before, and on 95\% of those days
the minute's lowest trade was at or below the 15:51 ask. A trade bar does not
show the size available, the queue, or what a particular order would have paid,
so the fills assumed here remain an approximation.

\subsection{Limit orders and late entry}
\label{sec:fills}

Since the trade made far more at the mid than at the ask, the obvious thing to
try is a limit order. I simulated six order styles, fixed before
any result was computed: paying the 15:51 ask (P0); a limit at the 15:51 mid for
one minute, else no trade (P1); the same, but counting a trade at the limit as a
fill (P1t); P1 followed by paying the 15:52 ask if unfilled (P2); a mid limit left
working for three minutes (P3); and a limit halfway between the mid and the ask,
then the 15:52 ask (P4). A limit counts as filled if a trade printed strictly
below it during a minute it was working, or if the quote at the end of that minute
offered the option at or below it. These are proxies: trade prints cannot show
whether a resting order would have executed, and the order is assumed to be live
with no delay.

\begin{table}[htbp]
\centering
\caption{Order styles for the rule, 2021--2026.}
\label{tab:fills}
\setlength{\tabcolsep}{4pt}
\begin{tabular}{lrrrrrr}
\toprule
& \multicolumn{4}{c}{2021--2026, 585 priced signal days} & \multicolumn{2}{c}{minus P0, \$/day} \\
\cmidrule(lr){2-5}\cmidrule(lr){6-7}
Order style & filled & \$/day & $t$ & minus P0 ($t$) & 2021--22 & 2023--26 \\
\midrule
P0 pay the 15:51 ask & 585 (100\%) & $+67.46$ & 2.29 & --- & --- & --- \\
P1 mid, 1 minute, else skip & 489 (84\%) & $+64.57$ & 2.37 & $-2.89$ ($-0.24$) & $-22.31$ & $+4.23$ \\
P1t as P1, touch counts & 501 (86\%) & $+61.04$ & 2.23 & $-6.42$ ($-0.56$) & $-22.31$ & $-0.59$ \\
P2 mid, 1 minute, then 15:52 ask & 585 (100\%) & $+69.63$ & 2.38 & $+2.17$ (0.55) & $+13.06$ & $-1.82$ \\
P3 mid, 3 minutes, else skip & 534 (91\%) & $+64.89$ & 2.31 & $-2.57$ ($-0.26$) & $+6.44$ & $-5.87$ \\
P4 halfway, 1 minute, then 15:52 ask & 585 (100\%) & $+62.23$ & 2.13 & $-5.23$ ($-1.63$) & $-2.36$ & $-6.29$ \\
\bottomrule
\end{tabular}
\tabnotes{\ITM{10}, \$0.09 fee, dollars per contract per priced signal day; an
unfilled order counts as \$0. The subperiods have 157 (2021--2022) and 428
(2023--2026) signal days. The $t$-statistics here treat signal days as
independent, so P0's $t$ is 2.29 rather than the per-day Newey--West 2.40 of
Table~\ref{tab:rpconf}. Descriptive.}
\end{table}

\begin{table}[htbp]
\centering
\caption{Why the mid limit does not help, and what a late entry costs.}
\label{tab:late}
\begin{tabular}{llrrr}
\multicolumn{5}{l}{\emph{Panel A. The same entries at the ask, split by whether P1's mid limit filled}} \\
\toprule
Period & P1 & days & \$/trade at the ask ($t$) & won \\
\midrule
2021--2026 & filled & 489 & $+42.52$ (1.31) & 52\% \\
 & \textbf{not filled} & \textbf{96} & $\mathbf{+194.49}$ \textbf{(2.80)} & \textbf{61\%} \\
\addlinespace
2013--2020 & filled & 69 & $+44.11$ (0.37) & 46\% \\
 & not filled & 147 & $-54.70$ ($-1.42$) & 40\% \\
\bottomrule
\end{tabular}

\bigskip
\begin{tabular}{lrrrr}
\multicolumn{5}{l}{\emph{Panel B. The same contract bought at a later minute's ask, \$/trade}} \\
\toprule
Buy at & 2013--2020 & 2021--2026 ($t$) & 2021--2022 & 2023--2026 \\
\midrule
15:51 & $-23.14$ & $+67.46$ (2.29) & $+115.46$ & $+49.85$ \\
15:52 & $-68.04$ & $+55.75$ (2.02) & $+111.07$ & $+35.46$ \\
15:53 & $-58.09$ & $+57.89$ (2.16) & $+128.30$ & $+32.06$ \\
15:55 & $-128.97$ & $+23.01$ (1.11) & $+51.10$ & $+12.71$ \\
\bottomrule
\end{tabular}
\tabnotes{The rule, \ITM{10}, \$0.09 fee, dollars per contract. Panel B uses the
days with a quote at every one of these minutes: 216 (2013--2020), 585
(2021--2026), 157 (2021--2022) and 428 (2023--2026), which is every priced day.
Descriptive.}
\end{table}

No order style beat paying the ask (Table~\ref{tab:fills}). On 2021--2026 every
style earned roughly what paying the ask earned, and every difference was
within noise. Panel A of Table~\ref{tab:late} suggests that the cause is
adverse selection. A limit at the mid fills when someone sells below the
mid, which happens most when the option is getting cheaper, that is, when the
trade is going wrong. When the trade works at once, the price moves away and the
limit is left unfilled. The days on which the mid limit filled earned well below
the average at the ask, and the days it missed earned far more. The limit saves
money on the days it fills and misses the best days, and the two effects roughly
cancel. Before 2021, with spreads about five times wider, the mid limit did help,
but on a small sample.

Entering late also costs money. Bought at a later minute's ask, the same
contract earned less the later the entry (Table~\ref{tab:late}, Panel B), about
\$12 a trade less at 15:52 and about two-thirds less at 15:55. This is relevant
because the results assume entry at 15:51:00 exactly, and a person reacting to
the candle would trade some seconds later.

\subsection{Selling before the close}
\label{sec:exits}

I also checked whether selling before the close would have done better. Each of
the 585 confirmation trades, bought
at the 15:51 ask, was also sold at the bid at the start of each later minute,
paying the fee twice. Every fixed exit minute from 15:52 to 15:59 earned less than
holding to settlement, by \$79 to \$136 a trade, and no year favoured an early
exit. The pre-registered model study of stop-loss and take-profit rules reached
the same conclusion (Section~\ref{sec:other}).

Rules that react to the option's price in real time did no better. A
pre-registered study tested rules that watch the option's own bid at each minute
from 15:52 to 15:58 and sell at the next minute's bid: 5 trailing stops (sell
when the bid falls 10--50\% below its peak since entry), 9 profit locks (once the
bid has risen 10--30\% above the price paid, sell if it gives back a third,
two-thirds or all of that gain) and 12 take-profits (sell once the bid is
10--30\% above the price paid, from 15:52, 15:54, 15:56 or 15:58). On the 270
trades of 2021--2023, a rule had to sell early on at least 25 trades and beat
holding in each half of the period; the best such rule would then have faced a
luck diagnostic and one test on 2024--2026. A first plan, which also included 18
rules that trail the index level implied by put--call parity, stopped at its
coverage gate before any rule was computed, because that level was often missing
in the last minutes. A revised plan, written after that stop, kept the 26
bid-based rules, whose quotes were complete.

No rule qualified, so nothing was tested. All 26 earned less than holding on
2021--2023, by \$9 to \$120 a trade. Take-profits did worst, \$49 to \$107 a
trade below holding, because they sell the trades that would have kept rising
and do nothing for the losers. I looked at 2024--2026 afterwards and it agrees:
take-profits trailed holding by \$39 to \$108 a trade, and the stops were within
\$31 of holding, none with $|t|$ above 1.6.

% =====================================================================
\section{Which candle?}
\label{sec:source}



Every signal so far was read from the CFD candle, but a trader would look at
the SPX index candle. The two track each other closely in level and direction, but a
candle's body-to-range ratio depends on its shape, and the shapes differ. On 36
sessions in the autumn of 2026 with both series, the index candle showed a strong
body at 15:50 on 20 days and the CFD candle on 14, all of them among the index
candle's 20. Over 2021--2026 the two ratios have a correlation of 0.78, and the
CFD's 0.60 corresponds to about 0.69 on the index candle. Two pre-registered
studies therefore read the 15:50 candle from archived one-minute bars, the SPX
index candle first and the SPY candle second, and traded the same SPXW option at
the prices of Section~\ref{sec:realprice}. The option outcomes of these days were
already known, so both studies are retrospective checks of the signal's source,
not independent confirmations.

The archives, both from CNBC, hold a valid 15:50 bar on 1{,}310 (index) and
1{,}312 (SPY) of the 1{,}322 sessions from 2021-01-01 to 2026-10-20 that have
option quotes and an official close. Checks made before each plan was locked, on
price levels only, confirmed the time stamps. For SPY they also showed that the
bars are traded prices, not dividend-adjusted, and that a second vendor's SPY
bars agreed to the cent on 14 of 22 sessions compared, with differences of up to
\$0.02 elsewhere. The signal is the same body test, computed in whole cents, and
the trade is unchanged: \ITM{10} in the candle's direction at the 15:51:00 ask,
held to settlement. Each primary line asks for a positive mean per day with
Newey--West $t \ge 2.0$; secondary lines also require a Holm-adjusted
$p \le 0.05$. The SPY study's second line was chosen after the agreement of two
candles had been seen to do well, and its lines are retrospective numerical
criteria.

\begin{table}[htbp]
\centering
\small
\renewcommand{\tabfont}{\fontsize{9}{11}\selectfont}
\caption{The rule read from three candle sources, 2021--2026.}
\label{tab:source}
\begin{minipage}{\linewidth}\centering
\textit{Panel A. The registered lines}\par\smallskip
\setlength{\tabcolsep}{4pt}
\begin{tabular}{@{}lrrrrrrrl@{}}
\toprule
Signal & Trades & \$/trade & At mid & Won & \$/day & $t$ & Holm $p$ & Reading \\
\midrule
CFD body $\ge 60\%$ (the rule) & 585 & $+\dollar67.46$ & $+\dollar105.90$ & 54\% & $+\dollar29.87$ & $+2.40$ &  & confirmed \\
\addlinespace
Index body $\ge 60\%$ & 710 & $+\dollar39.44$ & $+\dollar78.14$ & 52\% & $+\dollar21.39$ & $+1.53$ &  & not met \\
Index body $\ge 70\%$ & 579 & $+\dollar53.73$ & $+\dollar91.20$ & 53\% & $+\dollar23.75$ & $+1.86$ & 0.094 & insufficient \\
Index body 60--70\% & 131 & $-\dollar23.69$ & $+\dollar20.44$ & 49\% & $-\dollar2.37$ & $-0.41$ & 0.685 & insufficient \\
\addlinespace
SPY body $\ge 60\%$ & 523 & $+\dollar71.74$ & $+\dollar108.51$ & 53\% & $+\dollar28.62$ & $+2.35$ &  & met, retrospective \\
SPY and index both $\ge 60\%$ & 504 & $+\dollar80.44$ & $+\dollar117.34$ & 54\% & $+\dollar31.21$ & $+2.56$ & 0.0053 & met, retrospective \\
SPY body $\ge 70\%$ & 369 & $+\dollar103.40$ & $+\dollar137.67$ & 55\% & $+\dollar29.08$ & $+2.84$ & 0.0045 & met, retrospective \\
\bottomrule
\end{tabular}
\end{minipage}

\bigskip
\begin{minipage}{\linewidth}\centering
\textit{Panel B. Which of the three candles show a strong body}\par\smallskip
\begin{tabular}{@{}cccrrrrrr@{}}
\toprule
SPY & Index & CFD & Days & Trades & Ask & Mid & Won & $t$ \\
\midrule
yes & yes & yes & 436 & 435 & $+\dollar112.78$ & $+\dollar150.18$ & 56\% & $+3.11$ \\
yes & yes & -- & 69 & 69 & $-\dollar123.42$ & $-\dollar89.66$ & 36\% & $-1.73$ \\
yes & -- & yes & 9 & 9 & $-\dollar118.42$ & $-\dollar91.20$ & 33\% & $-0.80$ \\
yes & -- & -- & 6 & 6 & $-\dollar401.42$ & $-\dollar378.09$ & 33\% & $-1.25$ \\
-- & yes & yes & 82 & 82 & $-\dollar77.58$ & $-\dollar28.55$ & 44\% & $-1.20$ \\
-- & yes & -- & 121 & 121 & $-\dollar46.04$ & $-\dollar6.66$ & 50\% & $-0.82$ \\
-- & -- & yes & 54 & 54 & $-\dollar33.83$ & $-\dollar2.53$ & 54\% & $-0.49$ \\
-- & -- & -- & 523 & 0 & & & & \\
\bottomrule
\end{tabular}
\end{minipage}
\tabnotes{One \ITM{10} contract at the 15:51 ask, held to settlement, 2021-01-01
to 2026-10-20, dollars per contract. Panel A: the first row is the confirmation
of Table~\ref{tab:rpconf}; the index and SPY rows are the registered lines of the
two source studies, whose primary lines are the first of each block. Per day is
the mean of the daily series (0 on days outside the cell); $t$ is Newey--West,
lag 5; Holm $p$ is adjusted across each study's secondary lines. On the same days
as the index study the CFD rule made $+\dollar67.24$ a trade (583 trades), and on
the same days as the SPY study $+\dollar68.84$ (582). Panel B: \ITM{10} in the
direction of the first signalling candle in the order SPY, index, CFD; 1{,}300
sessions with all three candles; $t$ is a simple $t$-statistic; no session had
signalling candles in opposite directions. Panel B is descriptive.}
\end{table}

\paragraph{The index candle.} Read from the index candle, the rule made
$+\$39.44$ a trade, against $+\$67.24$ from the CFD candle on the same days, and
its daily series fell short of the criterion ($t = 1.53$; Table~\ref{tab:source},
Panel A). The band of 70\% or more came closer and also fell short, and the
60--70\% band showed no edge. The index candle is still informative about direction. After index candles of
70\% or more the index moved $+2.63$ basis points in the candle's direction ($t
= 4.55$), and after candles of 60--70\% it did not. The option trade falls
short because reading the index candle adds days without the effect. The 190 trades taken on days on which only the index candle showed a
strong body lost \$74.14 a trade, and 17\% of index candles of 70\% or more were
below 60\% on the CFD.

\paragraph{The SPY candle.} Read from the SPY candle, the rule made $+\$71.74$ a
trade ($t = 2.35$), about what the CFD candle made on the same days, and all three
registered lines met their retrospective criteria. The band of 70\% or more was
the strongest line and was positive in each of the six years; the 60--70\% band
made nothing, as on the other two sources. A strong SPY candle is nearly always a
strong CFD candle: 94\% of SPY candles of 70\% or more were CFD signals.

\paragraph{Three candles.} Panel B of Table~\ref{tab:source} shows the same
pattern across the three sources. The gain was confined to the 435 sessions on
which all three candles showed a strong body, and every other combination lost on
average, including the index and CFD candles signalling without SPY. No single
row other than the first is significant on its own, but together they suggest that a strong body on only some of the candles
carries no edge. One explanation, which I did not test, is that the CFD and SPY follow the futures, which move first,
while the index is computed from the last trades in 500 stocks and can show a
strong candle while it catches up. Because SPY was the third source tried on the
same outcomes, I treat these results as leads for a forward record and not as confirmations.

% =====================================================================
\section{The search for a better rule}
\label{sec:search}

The remaining studies looked for a better rule and did not find one.

\subsection{Re-tuning on recent data}
\label{sec:retune}

Tuning on 2013--2020 could not tell nearby settings apart, so a further
pre-registered study repeated the search over the same 2{,}178 settings on
2021--2023, the first three years of tight spreads and, from 2022, daily SPXW
expiries, and would test its pick once on 2024--2026. The pick had to lie
outside the rule's own family, pass the luck diagnostic and lead the rule on the
tuning years. Tuning and testing were separate runs, and the tuning step could
not read any date after 2023. Of 119 eligible settings, the best was the 15:50
candle with a body of at least 80\%, traded with the candle at \ITM{20}. It failed
the luck diagnostic ($p = 0.21$): about one shift in five found a best score at
least as good with signals unrelated to the days. No candidate advanced, and the rule was kept. Computed only after the decision
was recorded, the pick earned $+\$18.88$ a day on 2024--2026 against the rule's
$+\$29.97$, as one would expect of a winner chosen by luck.

I also could not tune the body threshold finely. In 2021--2026, the 15:50 candle
traded in its direction earned $-\$37$ a trade at the ask when its body was below
60\% of the range and $+\$67$ when it was at least 60\%, and the index move shows
the same break in 2013--2020. Above 60\%, per-day results are flat: the bootstrap
interval for the difference between thresholds of 0.70 and 0.60 runs from
$-\$10.7$ to $+\$8.8$ a day, and a walk-forward choice of threshold earned less
than a fixed 0.60. Within the signal, a larger ratio was not followed by a larger
gain. Across all days the ratio's correlation with what followed was significant
but faint, explaining at most 0.7\% of the day-to-day variation, and among signal
days it could not be distinguished from zero. The ratio therefore seems to work as a cut-off near 60\% and does not measure
the strength of the signal.

\subsection{Other closing minutes}
\label{sec:minutes}

A last pre-registered study asked whether another closing candle does better. It
covered 7{,}128 settings: three candle sources (the index candle, the SPY candle,
and the two agreeing in direction), the eleven one-minute candles from 15:45 to
15:55, twelve body cuts from 40\% to 95\%, trading with or against the candle,
and nine strikes from \ITM{20} to \OTM{20}. Each trade buys at the ask in the
snapshot that opens the next minute and holds to settlement. Settings were scored
on 2021--2023 as before, and 387 were eligible. The best one would be tested once
on 2024--2026 only if it was not a 15:50 setting, because those had already been
examined on the later years, and only if its luck diagnostic, repeated over 195
shifts, gave $p \le 0.05$. Before any statistic, the code reproduced the trades of
the two candle studies exactly. It also passed synthetic self-tests and a full
trial on synthetic data.

\begin{table}[htbp]
\centering
\small
\caption{The best eligible setting for each candle, 2021--2023.}
\label{tab:minutes}
\begin{tabular}{@{}llllrrrr@{}}
\toprule
Candle & Source & Body & Side & Strike & Score & \$/day & Trades \\
\midrule
15:45 & Index & $\ge 80\%$ & With & \ITM{20} & $+\dollar21.46$ & $+\dollar27.65$ & 166 \\
15:46 & Index & $\ge 40\%$ & With & \ITM{5} & $+\dollar1.48$ & $+\dollar6.81$ & 430 \\
15:47 & Index & $\ge 70\%$ & With & \ITM{5} & $-\dollar0.19$ & $+\dollar4.25$ & 235 \\
15:48 & Index & $\ge 40\%$ & With & \ATM & $+\dollar19.30$ & $+\dollar19.41$ & 416 \\
15:49 & Index & $\ge 90\%$ & With & \ATM & $+\dollar5.22$ & $+\dollar7.75$ & 104 \\
\textbf{15:50} & \textbf{Index} & $\mathbf{\ge 80\%}$ & \textbf{With} & \textbf{\ITM{20}} & $\mathbf{+\dollar37.46}$ & $\mathbf{+\dollar39.42}$ & \textbf{193} \\
15:51 & \multicolumn{4}{l}{No eligible setting} & & & \\
15:52 & Both & $\ge 50\%$ & Against & \ITM{5} & $+\dollar12.00$ & $+\dollar13.31$ & 237 \\
15:53 & Index & $\ge 80\%$ & With & \ATM & $+\dollar23.40$ & $+\dollar23.69$ & 177 \\
15:54 & Index & $\ge 60\%$ & With & \ITM{5} & $+\dollar33.55$ & $+\dollar32.43$ & 299 \\
15:55 & \multicolumn{4}{l}{No eligible setting} & & & \\
\bottomrule
\end{tabular}
\tabnotes{The highest-scoring eligible setting among those that read the given
candle, of 387 eligible among 7{,}128. Score is the mean profit per day averaged
with the same setting one cut above and below; \$/day is the setting's own mean
over 625 sessions. One contract at the ask in the snapshot that opens the next
minute, held to settlement, \$0.09 in fees. ``Both'' requires the index and SPY
candles to pass the cut in the same direction. Pre-registered.}
\end{table}

The search chose a 15:50 setting, the index candle with a body of 80\% or more
traded in its direction at \ITM{20}, so nothing went to the test. Its choice also
missed the luck bar: the shifted searches reached its score in 9 of 195 cases
($p = 0.051$), and its gain came almost entirely in the first half of the tuning
period. The search is still useful as a comparison of minutes. In
Table~\ref{tab:minutes} the 15:50 candle led, 15:54 came second, every other minute scored below both,
and two minutes had no eligible setting at all. Ranked afterwards over all of
2021--2026, all fifteen leading settings were 15:50 candles traded in their
direction, most of them at \ITM{20}; the first, both candles at 60\% or more,
made $+\$34.74$ a day ($t = 2.71$, 505 trades). That ranking reuses days already
seen and is not a test.

\subsection{Other studies}
\label{sec:other}

Four further studies, reported in full in the Internet Appendix, found nothing
that improved the rule.

\paragraph{Early exits.} Under the calibrated model, six stop-loss and
take-profit rules on the index path from 15:51 to 15:58 each earned less on
average than holding, at every strike and in every model variant. Stops cut the
worst loss by more than half but cost money on average. The likely explanation is that the signal is a small drift over nine minutes,
so many adverse moves reverse before the close, and every early sale pays the
spread again.

\paragraph{Intraday momentum and dealer gamma.} An adapted replication of the
market intraday momentum regression, on the CFD over 2013--2026, found a positive
but insignificant slope ($t = 1.02$), mildly positive only in 2013--2015 and about
zero on 2021--2026 ($t = -0.18$), the only period outside the samples of
\citet{ghlz2018} and \citet{bdlm2021}. This does not show that the published
results are wrong, since samples and instruments differ. No filter that kept the
candle trade only when it agreed with the day's earlier move improved it: trades
that agreed did better per trade, but the trades a filter skipped were still
profitable on average. A net-gamma proxy built from the open interest of every
SPX and SPXW option, assuming that dealers are long every call and short every
put, did not pass either test. Under model prices the candle trade did better on
negative-gamma days, but the difference missed its bar ($p = 0.029$ against
0.025) and was absent at real asks on 2021--2026.

\paragraph{Smaller contracts.} The SPX contract cost about \$1{,}100 on
2021--2026, so the same signal was traded on XSP (Mini-SPX) and SPY options about
a tenth the size. Neither was confirmed: XSP lost $\$49.66$ a trade ($t = -5.48$
per day) and SPY broke even. The SPX trade's edge was about six cents per dollar
of premium, and both smaller contracts cost more than that to trade; XSP's median
quoted spread was 31\% of the option's price, against 5\% for SPX on the same
days.

\paragraph{Forward evidence.} From 21 October to 6 November 2026, the strong-body
candle fired on seven of thirteen sessions and was right on six. Only four of those days have real one-minute option prices, which is too few to
conclude anything.

% =====================================================================
\section{Why 15:50?}
\label{sec:why}

Most of the evidence points to the same minute. The scan of 41 features
selected the 15:50 candle, and on 2021--2023 the re-tuning over 2{,}178
settings and the search over 7{,}128 both chose a 15:50 candle, although tuning
on 2013--2020, when spreads were wide
and only two settings were eligible, had favoured a 15:42 candle. A strong 15:50
candle is followed by an index move in its direction whether it is read from a
CFD, from the index or from SPY. I did not set out to explain the effect, but the timing of the closing auction
suggests a hypothesis. I describe it below and then report a pre-registered
test of its most direct prediction, which was not confirmed.

\subsection{The closing-auction hypothesis}
\label{sec:hypothesis}

From 15:50 the NYSE's regulatory closing imbalance is public and Nasdaq's early
imbalance indicator is disseminated, and the NYSE accepts new market-on-close and
limit-on-close orders only on the side that offsets the imbalance. The 15:50
one-minute candle is therefore the first full minute in which the market can
react to the published closing imbalance. A candle with a strong body is a minute
in which prices moved one way with little retracement, as might happen when a
large imbalance is published. If liquidity providers then absorb the imbalance
gradually up to the auction, the index would keep drifting in the same direction
for the remaining minutes. That is roughly the pattern in the data, a move of about two index points over
nine minutes, in the candle's direction, only when the candle has a
strong body. This reading fits the evidence of \citet{bm2023} that imbalance
information contributes to price discovery at the close, and it would explain why
the effect requires a one-sided minute and not just a large one.

Two facts do not fit the hypothesis easily. The signal was already present in
2013--2020, for most of which the NYSE's cut-off and mandatory publication came
five minutes earlier, at 15:45. Until late 2018 the only imbalance first published
at 15:50 was Nasdaq's, and for a few months after that no exchange began
publishing at 15:50, so the NYSE's schedule cannot fully explain it. Also, a strong body on the index candle alone, without the CFD and SPY candles, carried
no gain, which suggests that the futures-linked instruments react first. The exchanges' imbalance messages allow a direct test, reported next. I have
not tested whether the effect concentrates on days with large imbalances, or
whether it changed when the NYSE moved its publication from 15:45 to 15:50 in
April 2019, a few months after Nasdaq had moved its own from 15:50 to 15:55.

\subsection{A test with the published imbalances}
\label{sec:imbtest}

If traders react to the published imbalance during the 15:50 minute, the candle
should be right more often on days when it points the same way as that imbalance.
A pre-registered study (plan locked 2026-11-10 19:37 UTC, one registered run)
tested this prediction on the rule's 644 signal days of 2021--2026. For each year I took the 50 largest stocks in SPY at the previous year-end, from the fund's
N-PORT reports to the SEC (between 51\% and 63\% of its net assets), and bought
from Databento the imbalance messages that each stock's listing exchange sent
from 15:50. From each stock's first closing-imbalance record received in the
15:50 minute I summed the signed dollar imbalance, $D = \sum_i s_i Q_i p_i$, with
$s_i = +1$ for a buy imbalance, $-1$ for a sell imbalance and 0 for none, $Q_i$
the imbalance in shares and $p_i$ the reference price. A day agrees if the sign
of $D$ is the candle's direction. The primary test compares the rule's hit rate
on agreeing and disagreeing days, $d = h_a - h_d$, and passes if $d > 0$ and
$t \ge 2.0$. By the plan's power calculation it could detect a true $d$ of about
11 to 13 points, that is, only a large concentration of the edge on agreeing days. All 50 stocks had a usable record on every day, so every signal day
was classified.

\begin{table}[htbp]
\centering
\caption{Signal days that agree and disagree with the 15:50 closing imbalance, by year.}
\label{tab:imbyear}
\begin{tabular}{@{}lrrlrrlr@{}}
\toprule
 & & \multicolumn{3}{c}{Agreeing days} & \multicolumn{3}{c}{Disagreeing days} \\
\cmidrule(lr){3-5}\cmidrule(l){6-8}
Year & Days & $n$ & Hit rate (95\% CI) & bp & $n$ & Hit rate (95\% CI) & bp \\
\midrule
2021 & 111 & 92 & 62.0\% (52--71) & $4.67$ & 19 & 78.9\% (57--91) & $5.12$ \\
2022 & 105 & 76 & 57.9\% (47--68) & $4.00$ & 29 & 69.0\% (51--83) & $3.88$ \\
2023 & 113 & 73 & 53.4\% (42--64) & $0.33$ & 40 & 55.0\% (40--69) & $1.18$ \\
2024 & 106 & 69 & 59.4\% (48--70) & $2.59$ & 37 & 43.2\% (29--59) & $-1.90$ \\
2025 & 117 & 91 & 53.8\% (44--64) & $1.39$ & 26 & 69.2\% (50--83) & $3.97$ \\
2026 & 92 & 68 & 47.1\% (36--59) & $0.96$ & 24 & 70.8\% (51--85) & $4.29$ \\
\midrule
All & 644 & 469 & 55.9\% (51--60) & $2.41$ & 175 & 61.7\% (54--69) & $2.24$ \\
\bottomrule
\end{tabular}
\tabnotes{Signal days of the rule (a 15:50 body of 60\% or more), 2021--2026. A day agrees when the sign of the
summed dollar imbalance of the 50 largest stocks, from each stock's first closing-imbalance record received in the
15:50 minute, equals the candle's direction. Every signal day was classified, so no day is unclassified. Hit rate:
share of days on which the CFD moved from the 15:51 open to the 15:59 close in the candle's direction, with its
Wilson 95\% interval in percent; bp: the mean of that move in the candle's direction, in basis points. Over all
days the difference in hit rates is $-5.9$ points ($t = -1.35$), with a 20-day block-bootstrap 95\%
interval from $-14.8$ to $+4.0$ points.}
\end{table}

The candle and the imbalance usually pointed the same way, on 469 of the 644
days (72.8\%; Table~\ref{tab:imbyear}), but agreement did not make the candle
more accurate. It was right on 55.9\% of agreeing days and 61.7\% of
disagreeing days, $d = -5.9$ points ($t = -1.35$), with a block-bootstrap interval
from $-14.8$ to $+4.0$ points, and disagreeing days had the higher hit rate in
five of the six years, all but 2024. The imbalance direction alone was right on
51.1\% of days. The robustness checks fixed in the plan agree: the difference in
hit rates stayed negative with the latest record of the minute, with NYSE or
Nasdaq stocks alone and within year and direction, and the differences in the
average move and in the option P\&L were small against their noise ($t = 0.13$
and $0.79$; Internet Appendix). The prediction was therefore not confirmed.

The test has limited power, and a failure does not refute the auction
explanation. The price may adjust within the 15:50
candle itself, imbalances change after their first publication, and 50 stocks
are only part of the index. The high agreement is an association, not a cause, since the candle and the
imbalance can reflect the same closing orders. The test does show that the rule's edge was not concentrated on the days on which
its candle agreed with the imbalance first published at 15:50.

% =====================================================================
\section{Conclusion}
\label{sec:conclusion}

I asked whether a one-minute candle near the close predicts the
S\&P~500 into the close, and whether the prediction can be traded with same-day
options once the spread is paid. One candle passed a pre-registered, out-of-sample test, the 15:50 one-minute
candle with a strong body, traded in its own direction to the close. It is a
stricter version of an effect that was already known, and it has faded since
2023. Traded with a deep in-the-money 0DTE option bought at the quoted ask and
held to settlement, it made money on 2021--2026, though less than a pricing
model implied. The largest cost is the bid-ask spread, which took the whole
edge before 2021, and a limit order at the mid does not recover it because it
misses the best days. Re-tuned settings, real-time exit rules, momentum filters
and smaller contracts did not improve on the rule, and a search over the other closing candles chose the 15:50 candle again.

The evidence has limits. The edge is small, most of it comes from 2021--2022,
and it is fragile to execution, since an entry at 15:55 instead of 15:51 gave
up about two-thirds of it. The confirmation holds for the candle as
the backtests read it, from a CFD that tracks the index. Read from the SPX index candle under a plan locked in
advance, the rule earned less and did not meet its criterion. Read from the SPY candle
it met a retrospective one, and on every source the gain came on sessions
on which all three candles showed a strong body. A pre-registered test found that
the candle usually agrees with the closing-auction imbalance published at 15:50
but is no more accurate when it does, and because that test had limited power
the mechanism is still open. Other limits follow from the design. The holdout
was not blind, later studies reused the 2021--2026 outcomes, the fills are
proxies for what an order would have paid, and the luck diagnostic is
approximate.

One advantage of pre-registering every study is that the failed tests can be
reported and used. Because each plan was fixed and hashed before the data it
governs was read, a reader can see which refinements were tried and did not
help, and which explanation the data did not confirm. I read the evidence as a
small effect, sensitive to execution, that is worth continuing to measure but
is not an established edge.

There are three things I would do next. Settings suggested by the searches, such
as stricter 15:50 cuts and the agreement of several candles, should be tested only
on days after November 2026, under new registered plans. Closing-imbalance data
could test the rest of the hypothesis of Section~\ref{sec:hypothesis}, namely
whether the effect depends on the size of the imbalance and whether it changed
when the NYSE moved its publication to 15:50 in 2019. And a forward record of
actual order fills should replace the fill proxies, because only real fills can
show how the trade should be executed.

\paragraph{Data and code availability.} \ifblind The analysis code, the
pre-registration plans and their hash ledger are available on request (the link
is withheld for review). \else The analysis code, the pre-registration plans,
their hash ledger and the Internet Appendix are public at
\url{https://github.com/kv1514/spx-closing-candle-study}. \fi The OPRA data are
licensed from Databento and are not redistributed.

% =====================================================================
\begingroup
\setstretch{1.0}
\setlength{\bibsep}{6pt}
\begin{thebibliography}{99}

\bibitem[Bachelier(1900)]{bachelier1900}
Bachelier, L. (1900). Th\'eorie de la sp\'eculation. \textit{Annales
scientifiques de l'\'Ecole Normale Sup\'erieure}, 17:21--86.

\bibitem[Baltussen et al.(2021)]{bdlm2021}
Baltussen, G., Z. Da, S. Lammers, and M. Martens (2021). Hedging demand and
market intraday momentum. \textit{Journal of Financial Economics}, 142:377--403.

\bibitem[Beckmeyer et al.(2023)]{bbg2023}
Beckmeyer, H., N. Branger, and L. Gayda (2023). Retail traders love 0DTE
options\ldots\ but should they? SSRN Working Paper 4404704.

\bibitem[Bogousslavsky and Muravyev(2023)]{bm2023}
Bogousslavsky, V. and D. Muravyev (2023). Who trades at the close? Implications
for price discovery and liquidity. \textit{Journal of Financial Markets},
66:100852.

\bibitem[Dim et al.(2024)]{dev2024}
Dim, C., B. Eraker, and G. Vilkov (2024). 0DTEs: Trading, gamma risk and
volatility propagation. SSRN Working Paper 4692190.

\bibitem[Gao et al.(2018)]{ghlz2018}
Gao, L., Y. Han, S. Z. Li, and G. Zhou (2018). Market intraday momentum.
\textit{Journal of Financial Economics}, 129:394--414.

\bibitem[Harvey et al.(2016)]{harvey2016}
Harvey, C. R., Y. Liu, and H. Zhu (2016). \ldots\ and the cross-section of
expected returns. \textit{Review of Financial Studies}, 29:5--68.

\bibitem[Holm(1979)]{holm1979}
Holm, S. (1979). A simple sequentially rejective multiple test procedure.
\textit{Scandinavian Journal of Statistics}, 6:65--70.

\bibitem[Newey and West(1987)]{neweywest1987}
Newey, W. K. and K. D. West (1987). A simple, positive semi-definite,
heteroskedasticity and autocorrelation consistent covariance matrix.
\textit{Econometrica}, 55:703--708.

\bibitem[Nosek et al.(2018)]{nosek2018}
Nosek, B. A., C. R. Ebersole, A. C. DeHaven, and D. T. Mellor (2018). The
preregistration revolution. \textit{Proceedings of the National Academy of
Sciences}, 115:2600--2606.

\bibitem[Wilson(1927)]{wilson1927}
Wilson, E. B. (1927). Probable inference, the law of succession, and statistical
inference. \textit{Journal of the American Statistical Association}, 22:209--212.

\end{thebibliography}
\endgroup

\end{document}
